\documentclass[aspectratio=169]{beamer}

% -- Theme --------------------------------------------------------------------
\usetheme{default}
\usecolortheme{default}
\setbeamertemplate{navigation symbols}{}
\setbeamertemplate{footline}[frame number]
\setbeamertemplate{itemize items}[circle]

\definecolor{mainblue}{RGB}{0,70,127}
\setbeamercolor{frametitle}{fg=mainblue}
\setbeamercolor{title}{fg=mainblue}
\setbeamercolor{itemize item}{fg=mainblue}
\setbeamercolor{itemize subitem}{fg=mainblue}
\setbeamerfont{frametitle}{size=\large, series=\bfseries}

% -- Packages -----------------------------------------------------------------
\usepackage{booktabs}
\usepackage{amsmath}
\usepackage{xcolor}
\usepackage{colortbl}

\newcommand{\fhalf}{F\textsubscript{0.5}}
\newcommand{\errtag}[1]{\texttt{#1}}

% -- Metadata -----------------------------------------------------------------
\title{Multi-Class Grammatical Error Detection for Correction:\\
A Tale of Two Systems}
\author{Authors: Yuan, Taslimipoor, Davis \& Bryant \\ \small University of Cambridge}
\date{Paper from: EMNLP 2021}

% -----------------------------------------------------------------------------
\begin{document}
% -----------------------------------------------------------------------------

\begin{frame}
  \titlepage
\end{frame}

% -- MOTIVATION ---------------------------------------------------------------

\begin{frame}{The Problem}
  GEC systems correct errors but are blind to what \textit{kind} of error
  they are dealing with.\\[1em]

  GED (grammatical error detection) traditionally just flags tokens as
  correct or incorrect -- a binary signal that throws away type information.\\[1em]

  \textbf{Can richer, multi-class error type information improve GEC?}
\end{frame}

\begin{frame}{The Idea}
  Train a GED system that predicts fine-grained error types,
  not just correct vs.\ incorrect.\\[1.5em]

  Feed those predictions into a GEC system in two ways:

  \begin{itemize}
    \item During \textbf{fine-tuning}: give the GEC model error type labels
      as additional input
    \item During \textbf{inference}: use GED predictions to re-rank
      the GEC system's candidate outputs
  \end{itemize}

  \vspace{1em}
  The two methods turn out to be complementary and are combined
  in the final system.
\end{frame}

% -- GED SYSTEM ---------------------------------------------------------------

\begin{frame}{The GED System: ELECTRA}
  Three pre-trained models tested for binary GED: BERT, XLNet, ELECTRA.\\[1em]

  ELECTRA performs best. Its pre-training objective -- distinguishing real
  tokens from plausible replacements -- is conceptually close to GED.\\[1em]

  Each model is fine-tuned with a linear classification layer on top,
  treating GED as a token-level sequence labelling task.

  \vspace{1em}
  \begin{center}
  \begin{tabular}{lcc}
    \toprule
    \textbf{Model} & \textbf{BEA-dev \fhalf} & \textbf{FCE-test \fhalf} \\
    \midrule
    Bell et al.\ (2019) -- prior SOTA & 48.50 & 57.28 \\
    GED (BERT) & 59.23 & 67.88 \\
    GED (XLNet) & 63.23 & 69.75 \\
    \textbf{GED (ELECTRA)} & \textbf{65.54} & \textbf{72.93} \\
    \bottomrule
  \end{tabular}
  \end{center}
\end{frame}

\begin{frame}{Multi-Class GED: Four Levels of Granularity}
  Using ERRANT's modular tag structure, ELECTRA is extended to predict
  error types at three levels beyond binary:

  \vspace{1em}
  \begin{center}
  \begin{tabular}{lll}
    \toprule
    \textbf{Setting} & \textbf{Classes} & \textbf{Example label} \\
    \midrule
    Binary   & 2  & \errtag{I} (incorrect) \\
    4-class  & 4  & \errtag{M} (missing) \\
    25-class & 25 & \errtag{VERB:TENSE} \\
    55-class & 55 & \errtag{R:VERB:TENSE} \\
    \bottomrule
  \end{tabular}
  \end{center}

  \vspace{1em}
  Key finding: adding more classes barely affects \textit{binarised} detection
  performance. The models detect roughly the same errors regardless of
  granularity -- they just struggle to classify the specific type.
\end{frame}

% -- MULTI-ENCODER GEC --------------------------------------------------------

\begin{frame}{Using GED as Auxiliary Input: Multi-Encoder GEC}
  A second encoder is added to the standard Transformer encoder-decoder
  GEC model. It processes GED predictions alongside the source sentence.\\[1em]

  A gating mechanism blends the two encoder representations:
  \[
    \text{Gating}(\text{MH}) = \lambda \cdot \text{MH}_\text{src}
    + (1 - \lambda) \cdot \text{MH}_\text{ged}
  \]
  where $\lambda$ is learned from the data.\\[1em]

  \textbf{Two-step training:}
  \begin{itemize}
    \item Step 1: train the base GEC model on full parallel data as normal
    \item Step 2: freeze the base model, train only the new GED components
      on a smaller high-quality dataset with GED labels added
  \end{itemize}

  Freezing the base model prevents overfitting when the fine-tuning
  set is much smaller than the pre-training set.
\end{frame}

\begin{frame}{Fine-Tuning Results: 4-Class Works Best}
  Adding GED predictions as auxiliary input consistently improves over
  the baseline. The 4-class system performs best.\\[1em]

  \begin{center}
  \begin{tabular}{lc}
    \toprule
    \textbf{System} & \textbf{BEA-dev \fhalf} \\
    \midrule
    Baseline (no GED) & 49.80 \\
    \midrule
    Binary GED & 52.20 \\
    \textbf{4-class GED} & \textbf{52.84} \\
    25-class GED & 51.78 \\
    55-class GED & 51.52 \\
    \midrule
    55-class Oracle (upper bound) & 70.24 \\
    \bottomrule
  \end{tabular}
  \end{center}

  \vspace{0.5em}
  4-class wins because it strikes the best balance between
  label informativeness and prediction reliability.
  The large oracle gap shows the bottleneck is GED accuracy, not model capacity.
\end{frame}

% -- RE-RANKING ---------------------------------------------------------------

\begin{frame}{Using GED for Re-Ranking}
  At inference time, the GEC decoder produces a 10-best list of candidate
  corrections. GED is used to pick the best one.\\[1.5em]

  \textbf{Procedure:}
  \begin{itemize}
    \item Convert each candidate correction to token-level error labels
      using ERRANT
    \item Run the GED system on the source sentence to get predicted labels
    \item Pick the hypothesis whose labels best agree with the GED predictions
      (minimum Hamming distance)
  \end{itemize}

  \vspace{1em}
  No GEC system scores or other features are used -- only GED predictions.
  The method can be applied to output from any GEC system.
\end{frame}

\begin{frame}{Re-Ranking Results: 55-Class Works Best}
  For re-ranking, more granular GED labels are better.
  55-class outperforms all coarser alternatives.\\[1em]

  \begin{center}
  \begin{tabular}{lcc}
    \toprule
    \textbf{GED for re-ranking} & \textbf{BEA-dev \fhalf} \\
    \midrule
    Baseline (no re-ranking) & 49.80 \\
    \midrule
    Binary & 51.21 \\
    4-class & 52.42 \\
    25-class & 52.68 \\
    \textbf{55-class} & \textbf{54.35} \\
    \bottomrule
  \end{tabular}
  \end{center}

  \vspace{0.5em}
  55-class labels are more discriminative between candidates,
  even if noisier. For re-ranking, discrimination matters more
  than reliability.
\end{frame}

\begin{frame}{Why Different Granularities Win in Different Roles}
  The right level of detail depends on how the GED information is used.\\[1.5em]

  \textbf{Fine-tuning (4-class wins):}
  \begin{itemize}
    \item Labels are baked into model weights during training
    \item Noisy fine-grained labels corrupt the training signal
    \item Coarser but reliable labels are more useful
  \end{itemize}

  \vspace{1em}
  \textbf{Re-ranking (55-class wins):}
  \begin{itemize}
    \item Labels are used to compare candidate outputs at inference time
    \item Fine-grained labels better discriminate between similar hypotheses
    \item Noise matters less when you are choosing between options
  \end{itemize}
\end{frame}

% -- FINAL RESULTS ------------------------------------------------------------

\begin{frame}{Final System and State of the Art}
  Combining 4-class fine-tuning and 55-class re-ranking gives further gains,
  confirming the two methods are complementary.\\[1em]

  \begin{center}
  \begin{tabular}{lcc}
    \toprule
    \textbf{System (single model, unconstrained)} & \textbf{BEA-test \fhalf} & \textbf{CoNLL-14 \fhalf} \\
    \midrule
    Stahlberg \& Kumar (2021) & 70.4 & 66.6 \\
    Kaneko et al.\ (2020) + GED & 65.6 & 62.6 \\
    \midrule
    Baseline (this work) & 65.1 & 59.4 \\
    + Multi-encoder GEC & 67.6 & 62.7 \\
    \textbf{+ GED re-ranking} & \textbf{70.6} & \textbf{63.5} \\
    \bottomrule
  \end{tabular}
  \end{center}

  \vspace{0.5em}
  New single-model NMT SOTA on BEA-test. Outperforms all prior systems
  that combine GED and GEC.
\end{frame}

% -- CONCLUSION ---------------------------------------------------------------

\begin{frame}{Summary}
  \begin{itemize}
    \item ELECTRA fine-tuned for GED sets a new binary detection SOTA and
      extends naturally to multi-class\\[0.8em]
    \item A multi-encoder GEC model with gated attention incorporates GED
      predictions during fine-tuning\\[0.8em]
    \item GED predictions also improve GEC as a post-hoc re-ranking signal\\[0.8em]
    \item The two methods prefer different granularities: 4-class for
      fine-tuning, 55-class for re-ranking\\[0.8em]
    \item Oracle experiments show large headroom: the bottleneck is GED
      accuracy, not model capacity\\[0.8em]
    \item Combined system sets a new NMT SOTA on BEA-test
  \end{itemize}
\end{frame}

\end{document}
